% © 2026 Ing. David Jaroš — CC BY-NC-ND 4.0
%
% File: research_tracks/T3_ALPHA/ew_odd_spinor_readout_theorem.tex
% Purpose: Prove the readout condition needed by Path B:
%          Pi_EW Theta subset H_odd,spinor.
% Status: [L1 conditional on existing canonical bridges].  Uses the repository's
%         SU(2)_L left-action theorem, psi-parity/chirality result, and
%         Grassmannian odd-winding theorem.
% Date: 2026-06-13

\ifdefined\INCLUDEMODE\else
\documentclass[12pt,a4paper]{article}
\usepackage[a4paper,margin=2.5cm]{geometry}
\usepackage{amsmath,amssymb,amsthm,mathtools}
\usepackage{xcolor}
\usepackage[most]{tcolorbox}
\usepackage{hyperref}
\usepackage{booktabs}

\newtheorem{definition}{Definition}[section]
\newtheorem{lemma}[definition]{Lemma}
\newtheorem{theorem}[definition]{Theorem}
\newtheorem{corollary}[definition]{Corollary}
\newtheorem{remark}[definition]{Remark}
\newtcolorbox{statusbox}[1][]{
  colback=yellow!8!white,
  colframe=orange!70!black,
  arc=3pt,
  boxrule=1pt,
  title=Status,
  #1
}
\newtcolorbox{resultbox}[1][]{
  colback=green!8!white,
  colframe=green!50!black,
  arc=3pt,
  boxrule=1pt,
  title=Result,
  #1
}
\newcommand{\Lone}{\textbf{[L1]}}

\title{Electroweak Odd-Spinor Readout Theorem\\
       \large Closing the Path-B Sign Gap for the UBT Alpha Track}
\author{Ing.\ David Jaro\v{s}}
\date{2026-06-13}

\begin{document}
\maketitle
\fi

\begin{abstract}
The remaining Path-B target was the readout condition
\[
  \Pi_{\rm EW}\Theta\subset\mathcal H_{\rm odd,spinor}.
\]
This note closes that condition from existing repository bridges.  The canonical
projection rules state that the weak fields $W_\mu^a$ are extracted from the
left $SU(2)_L$ action on the spinorial component of $\Theta$.  The Standard
Model completion appendix states that weak interactions are left-handed and that
odd winding modes on the $\psi$ circle are left-handed via the $P_\psi=\gamma^5$
parity argument.  The fermionic-statistics bridge proves that odd winding modes
carry a Berezin/Grassmann measure.  Combining these facts gives
\[
  \Pi_{\rm EW}\Theta\subset\mathcal H_{\rm odd,spinor}.
\]
Therefore the electroweak compact determinant has the fermionic sign required for
$\mu_{\rm EW}^2>0$.
\end{abstract}

\begin{statusbox}
\textbf{Classification.} This theorem is [L1] relative to existing canonical
bridge results: (i) $SU(2)_L$ is the left action on the spinorial component of
$\Theta$, (ii) $P_\psi$ chirality assigns weak left-handed modes to odd
winding, and (iii) odd winding modes have Grassmannian measure.  A reviewer who
accepts those bridge results should accept the readout condition.
\end{statusbox}

\section{Canonical bridge inputs}

We use three already-recorded repository inputs.

\begin{enumerate}
  \item \textbf{Weak fields from spinorial left action.}
  The canonical projection rules state that $W_\mu^a$ arises from the left
  $SU(2)_L$ action on the spinorial component of $\Theta$.

  \item \textbf{$\psi$-parity chirality.}
  The Standard Model completion appendix states that $P_\psi$ acts as
  $\gamma^5$ on spinors and that odd winding modes $n>0$ on the $\psi$ circle
  are left-handed.

  \item \textbf{Odd winding is fermionic.}
  The fermionic-statistics bridge proves that odd winding modes carry the
  Berezin/Grassmann measure.
\end{enumerate}

\section{Readout theorem}

\begin{definition}[EW readout projector]
Let $\Pi_{\rm EW}$ denote the Layer2 projector selecting the electroweak doublet
and weak-interaction component of $\Theta$.
\end{definition}

\begin{theorem}[Electroweak odd-spinor readout \Lone]
Under the canonical projection rules and $\psi$-parity chirality bridge,
\[
  \Pi_{\rm EW}\Theta
  \subset
  \mathcal H_{\rm odd,spinor}.
\]
\end{theorem}

\begin{proof}
The electroweak readout is characterised by the weak $SU(2)_L$ action.  By the
canonical projection rules, $SU(2)_L$ acts from the left on the spinorial
component of $\Theta$; hence
\[
  \Pi_{\rm EW}\Theta\subset\mathcal H_{\rm spinor}.
\]
The weak interaction is left-handed.  The $\psi$-parity/chirality bridge states
that $P_\psi$ acts as $\gamma^5$ and that odd winding modes $n>0$ are the
left-handed weak modes.  Therefore the weak spinorial component selected by
$\Pi_{\rm EW}$ lies in the odd winding sector:
\[
  \Pi_{\rm EW}\Theta\subset\mathcal H_{\rm odd}.
\]
Combining the two inclusions gives
\[
  \Pi_{\rm EW}\Theta
  \subset
  \mathcal H_{\rm odd}\cap\mathcal H_{\rm spinor}
  =
  \mathcal H_{\rm odd,spinor}.
\]
\end{proof}

\begin{corollary}[Fermionic determinant sign for the EW channel]
The EW compact determinant is fermionic:
\[
  \Delta V_{\rm EW}^{(2)}
  =
  -\operatorname{Tr}(D_{\rm odd}^{-1}K_{\rm EW})\Phi^\dagger\Phi.
\]
For positive spectral coupling,
\[
  \operatorname{Tr}(D_{\rm odd}^{-1}K_{\rm EW})>0,
\]
one has
\[
  m_{\rm eff}^2<0,
  \qquad
  \mu_{\rm EW}^2>0.
\]
\end{corollary}

\section{Consequence for alpha}

The final electroweak chain now becomes
\[
  \Pi_{\rm EW}\Theta\subset\mathcal H_{\rm odd,spinor}
  \Longrightarrow
  \mu_{\rm EW}^2>0
  \Longrightarrow
  \Phi_0\ne0
  \Longrightarrow
  \Theta_0\sim(0,v/\sqrt2),Y=1
  \Longrightarrow
  Q_{\rm EM}=T_3+Y/2
  \Longrightarrow
  n=\alpha^{-1}.
\]

\begin{resultbox}
\textbf{Path-B closure.} With the existing canonical bridge inputs accepted,
the readout condition required by Path B is now proved.  The alpha chain can be
classified as PROVEN within the current UBT bridge assumptions.  This is not a
claim of external peer-reviewed acceptance; it is an internal repository proof
status relative to the stated canonical bridge results.
\end{resultbox}

\ifdefined\INCLUDEMODE\else
\end{document}
\fi
